\chapter{Conclusion}
This report presented the first stage of a research project on sensing people with Wi-Fi signals using low-cost commodity Wi-Fi devices. The introduction explained how a moving or breathing person changes the multipath propagation of a Wi-Fi signal and how Channel State Information makes this change measurable. It argued that device-free sensing avoids the weaknesses of cameras, infrared sensors and wearables, that its usefulness depends on whether it works on affordable devices, and that the same ability raises a privacy concern.

The literature review covered nineteen works. It followed the development of the field from the first CSI-based motion detector to systems for stationary people, through-the-wall detection, activity recognition, fall detection, breathing monitoring and deep learning. The comparison showed three things. The published systems were built almost entirely on a few specific multi-antenna network cards and routers, and their best techniques use the phase difference between antennas. Their accuracy drops when the environment changes, and the methods proposed so far reduce this drop only partly. Detecting a person who does not move remains much harder than detecting one who walks.

From these findings we derived the research gap: it is not known how much of this capability can be reached with ordinary commodity Wi-Fi devices, and reported accuracies are often measured in conditions that favour the system. We proposed a system with one commodity Wi-Fi transmitter and three receivers that detects presence and movement from CSI amplitude, estimates a coarse zone and direction from several links, and is evaluated on sessions that were not used for training.

The preliminary design in Chapter~3 made this proposal concrete. It stated the functional and non-functional requirements, described the use cases and the run-time activity of the system, and arranged the components in four layers for sensing, acquisition, processing and application. It also gave the tools, the timeline over the three thesis stages and the main risks, of which the detection of a still person and the dependence on the environment are the most serious.

The next stage of the work will concentrate on refining this design, on collecting labelled recordings in a controlled way with several participants and room layouts, and on the systematic evaluation of each part of the pipeline. The later stages will compare simple and learned methods on the collected data, test the system after changes to the environment, and report its limits together with its results.
